\section[Capítulo 2: Aprendizaje Práctico de Ciberseguridad]{\centering Capítulo 2: Aprendizaje Práctico de Ciberseguridad}

El aprendizaje práctico es necesario porque la ciberseguridad no se limita a saber conceptos, sino a usarlos en momentos donde la información puede ser incompleta o ambigua. Un estudiante puede conocer qué es el \textit{phishing}, pero aun así fallar si no sabe revisar dominios, remitentes o señales de urgencia falsa (Prince, 2004).

Este capítulo explica por qué el proyecto se apoya en actividades interactivas. La propuesta no busca reemplazar la clase teórica, sino complementarla con ejercicios donde el estudiante tome decisiones, observe consecuencias y refuerce su criterio de respuesta (Kolb, 1984).

\subsection{Limitaciones del enfoque teórico}

La enseñanza expositiva permite introducir conceptos, pero no siempre alcanza para desarrollar habilidades de análisis y respuesta. En seguridad informática, muchas tareas requieren práctica: revisar un correo, interpretar una alerta, decidir si se escala un caso o contener un equipo comprometido (Prince, 2004).

En grupos numerosos, el docente puede tener dificultades para acompañar a cada estudiante durante ejercicios prácticos. Esto reduce oportunidades de retroalimentación y puede dejar una brecha entre lo que el estudiante sabe explicar y lo que realmente puede resolver frente a un caso (Biggs \& Tang, 2011).

\subsection{Aprendizaje activo}

El aprendizaje activo plantea que el estudiante participa mediante análisis, resolución de problemas y toma de decisiones. Este enfoque encaja con ciberseguridad porque obliga a usar evidencias, evaluar alternativas y justificar acciones frente a una situación concreta (Prince, 2004).

\begin{longtable}{|p{0.30\textwidth}|p{0.55\textwidth}|}
\hline
\textbf{Estrategia} & \textbf{Aplicación en el videojuego} \\
\hline
Análisis de casos & Correos, llamadas y alertas con información incompleta o sospechosa (Kolb, 1984). \\
\hline
Toma de decisiones & Acciones como aprobar, escalar, reportar o contener un incidente (Prince, 2004). \\
\hline
Aprendizaje por error & Consecuencias mostradas en reportes de turno y registros del sistema (Hattie \& Timperley, 2007). \\
\hline
Progresión gradual & Casos simples al inicio y escenarios combinados en niveles posteriores (Biggs \& Tang, 2011). \\
\hline
\end{longtable}

Desde esta perspectiva, el videojuego funciona como un espacio de práctica controlada. El estudiante puede ensayar respuestas y equivocarse sin generar daño real, lo cual resulta útil para temas donde el error en un entorno laboral puede tener consecuencias serias (Kolb, 1984).

\subsection{Retroalimentación en el aprendizaje}

La retroalimentación ayuda a reducir la distancia entre el desempeño actual y el desempeño esperado. Para que sea útil, debe entregar información sobre la acción realizada y orientar una mejora, no solo indicar si algo estuvo bien o mal (Hattie \& Timperley, 2007).

En el proyecto, la retroalimentación se plantea mediante consecuencias dentro del escenario. Por ejemplo, si el jugador aprueba un correo sospechoso, el reporte posterior puede indicar que un equipo inició actividad anómala. Así se mantiene el tono de simulación y, al mismo tiempo, el estudiante entiende el efecto de su decisión (Hattie \& Timperley, 2007).

\subsection{Síntesis del capítulo}

El aprendizaje práctico justifica el uso de escenarios interactivos en el proyecto. Para formar habilidades de ciberseguridad, el estudiante necesita observar, decidir y revisar consecuencias. El videojuego se plantea como una forma de apoyar ese proceso sin abandonar el contenido académico (Prince, 2004).

\newpage
